\section{NeoPlatonism \& Orthodoxy}

An Eastern Orthodox response\footnote{\url{http://orthodoxinfo.com/inquirers/orth\_plato.aspx}} to Protestant critique reveals some very interesting facts about the tradition.

First of all, it considered itself to \emph{build} upon classical tradition:

\begin{quotex}
Towards late antiquity, two words were in philosophical use which meant substance or substantive being. As Aristotle put it, being can be said in many ways or there are many meanings of Being. He isolated and identified Being according to the different categories of which substance (ousia) was the most fundamental, Being as act and potency, Being as true, and Being as contingent as opposed to necessary/essential. There has been on and off again the debate over the question of whether this listing of the meanings of Being is complete. We will return to that in a bit. The Byzantine answer was it was incomplete.
\end{quotex}


There is no difficulty here, if it is pointed out that often Christianity did not do this; the intent is what interests us, not the result (by definition, imperfect, since it is ``realized'' on the exoteric plane). Byzantium intended to perfect the revelation given to the Greeks and Romans, not to supplant it.

The next, longer passage is very relevant for Calvinists and their descendents, the hyper-Calvinist materialists called ``Progressives'':

\begin{quotex}
The present is the predestined creature of the past both in Calvin and Newton. The essentialism downgrades accidental being or sumbebekos. Now \emph{sumbebekos} as accidental being was also not well understood in the west. It is not being that exists in another like a form of being found within one of the predicate categories. It designates the happenstantial, the contingently situational so it is not precisely necessary-essential that my wife has auburn hair; that she has this hair color is accidental in the sense of \emph{sumbebekos}, it is fortuitous in the sense that her hair could have indifferently been auburn or not auburn but it is necessary that it have some color. This kind of accidental being can be found in every category. Thus, in the pagan Greek view, that a human is this or that individual is accidental but it is necessary that a human be somebody. The distinction then between accidental being in the sense of s\emph{umbebekos} from an accident-predicate is that the accidental being is the contingent or fortuitious actuality according to every necessary category. So, accidental being in this sense is the contingent affection or actual constellation or situation or event that takes place according to the necessity of different categories. On the Orthodox reading, this \emph{sumbebekos} is really God's governance or the concrete realization of God's purpose in time. To downgrade it is a spiritually sick symptom of the religiosity of \emph{sarx} that finds time to be a corrosive force that bears no faith or hope. Thus, accidental being just as the unique particularity of persons and the future, is downgraded in pagan metaphysics in favor of fully determinate being-essences (the dead, fully completed and determinate and deterministic past idealized as enduring presence) or ousiological being for the paradigmatic past participle. It was not till Scotus developed his concept of \emph{haecceity} and the individual's will did you find the beginning of the overturn of a ousiological metaphysics in western Christian theology. Maximos pointed out the original meaning of sumbebekos was not negative as it as in pagan philosophical theology but was positive. It meant the adventurous and creative way things are ``brought together.''
\end{quotex}

This summarizes how Eastern Churches have been able retain their ``ethnicity'' and ``race'' -- the Serbian race is incidental to Grace (in one sense, it is recognized), but according to Providence, such ``incidents'' are not dispensable from the essence -- one cannot ``extract'' the essence of Eastern Orthodoxy apart from the Slavic soul. Of course, this is precisely the goal of esoteric studies -- to draw the ``soul'' (and even the body) upward into the Spirit in order to ``save'' them. The spirit enters heavens bearing that which it has conquered, bringing trophies. By contrast, the Western Church has fallen afoul of a very literal reading of both Aquinas \& Aristotle, in which abstract categories of mere intellection substitute for the living organic reality of an actual, suffering, marching Church.

\begin{quotex}
In this, emphasize that Trinity and Scripture requires a strong ecclesiology---a weakness in post-Reformed Protestant thought which is docetic in its ecclesiology---because the primary framework of the agape relation to neighbor is Church-building either through missions to bring them in or edification to build them up within the Body of Christ. In the original sources, even Luther does not use the construction ``the Christian religion'' or ``Christianity is a religion'' even though you will find ``the religion (synonymous with piety) of Christians'' because even for him Christianity = Church. It is Calvin where this equation is loosened, and thus, ecclesiology becomes very weakened.
\end{quotex}

The Church is the Church. There is no ``Christianity'', only the Church. As such, Gornahoor aims to strengthen the older ideal of the Western Church. Not the concept, but the ideal, which is very different. Calvinism has some promise in its tragic doctrine of depravity, yet a universe in which God ``restlessly orders every literal event as if by an unseen hand perhaps the free market? '' is not a doctrine that will save the West.

\begin{quotex}
The Orthodox understanding is that the pagan concept of divine immutability was an allergic reaction to time and contingency that needed to find a time-proof foundation of necessity for religious reassurance. Time was the evil Kronos who ate its children. So, the concept is symptomatic of a spiritual sickness of the religiosity of sarx (which in Orthodoxy is not ``flesh,'' but deathliness, related to sarcophagus, it is the opposite of spirit, it is de-spiritedness that emotionally believes in death more than God despite what the head believes). Pagan philosophical theology sought a predestinied necessity in a time-proof order or an anti-time God. Since everything human is temporal, one gets the unknown God because God is what anything temporal-mortal is not.
\end{quotex}


Despite the fact that this is an un-nuanced view of Classical antiquity, it is at least promising in that it attempts to correct and build upon it organically, rather than unconsciously assimilating it through denial.

However, it appears the author is willing to understand parts of antiquity on its own:

\begin{quotationx}
Into this context comes the tradition of Judaism. Philo uses these distinctions to suggest that the personal God of Israel is a uniquely distinctive and singular reality or hypostasis in contrast to the impersonal thought thinking itself thing (ousia) of Aristotle. One branch of Middle Platonism picks this up and God and souls become \emph{hypostases} while impersonal things, including cups, chairs, etc., become \emph{ousias}. The other branch of Middle Platonism represented by Numenius of Apamea pushes the God as impersonal monad and ousia interpretation which nevertheless reinforces the \emph{ousia} = non-personal substance/thing of a typical kind and \emph{hypostasis} = personal substance of a uniquely and distinctive individual (usually God as person).

We have here the elements of a conceptual revolution in ancient philosophy and religion. It was the characteristic trait of ancient philosophy, religion, and humanism that the individual was not valued in and of itself in its particularity but only as an paradigmatic instance of a universal ideal type. Person was just an epiphenomenal mask. By contrast, as Tillich brings out, even the post-Christian humanism of the modern world is Judeo-Christian to the extent the unrepeatable individual \emph{per se} is valued in and of itself. The roots of this contrast is in the Trinitarian controversy and the debt the parties owed to Philo and Origen.
\end{quotationx}

This shows Judaism's role in building Tradition (via Philo), while also nuancing various branches of Platonism (eg., Syriac). I don't see how the following sentence doesn't really summarize the basic understanding of esoteric Christianity:

\begin{quotex}
Human existence is comprised of human nature (\emph{ousia}), with its matter and form, as a species principle of nature (\emph{logos physeos}) that is individuated by the personal mode of existence (\emph{tropos hyparxeos}) of the hypostasis. But Orthodoxy has a developmental view. Even if there had been no fall, Adam would have still had to evolve from the Image into the Likeness through the Incarnation (anselmian theory is wrong on the eastern view if purpose of Incarnation was just to correct the fall because Church is the goal as the Body of Christ) in a process called theosis (deification) in synergy with the divine Energies. While in principle it is the hypostasis that individuates, instead of one of the metaphysical elements of the human ousia, it is the process of theosis that is the actual process of individuating the one becoming a person in the fulness of the Likeness of the interpersonality of the Trinitarian life...
\end{quotex}

Eastern Orthodoxy deserves to be studied by esotericists, who should begin with Mouravieff's \emph{Gnosis}.

Here is an example of how a concrete Christian communion should go about incorporating and developing ancient teachings; it also instructs us how abstract theological debates \emph{were part \& parcel of the ancient thought-world}, and not ``inventions'' of the theologians. They were discussing (at Nicea) what the Christ-even meant to them, in terms of what they already believed.


\flrightit{Posted on 2011-09-17 by Logres}

\begin{center}* * *\end{center}

\begin{footnotesize}
\begin{sffamily}
\texttt{Bob on 2011-09-17 at 10:54 said:}

First, let me say that I am a neophyte to all things esoteric and Traditionalist.

I have recently been introduced to Evola and Guenon, and have found it somewhat strange that they barely ever mention the Eastern Orthodox Church when considering the merits of Christianity. There is always a dichotomy between East and West in Guenon's writings, but the ``East'' seems to refer to Buddhism, Hinduism, Islam etc.

I'm currently reading the writings of Vladimir Lossky, and I'm somewhat startled at how much mysticism and Hellenic thought is present within Orthodoxy.

It seems clear to me that, at the very least, Orthodoxy has been fairly successful in warding off modern influence and preserving a lot of ancient tradition and traditional thought. As such, I find myself drawn to it.

Could Eastern Orthodoxy (in tandem with the re-emergence of Russia) play a part in a revitalization of the West?

\hfill

\texttt{Logres on 2011-09-17 at 16:26 said:}

Probably so, and you'll want to read Matthew Raphael Johnson, which Cologero has linked to at this site: \url{http://www.gornahoor.net/?p=2842}

He has some good articles on the web: 
\url{https://thattimehascome.blogspot.com/2011/05/frequently-asked-questions-about.html}


I've toyed with the idea of using ``Western rite Orthodoxy'' as a vehicle for what is being discussed here. The Western rite is being treated like a red-headed stepchild anyway, inside the Orthodox Church (from what I gather) so no one should care. It isn't part of the ``Church'' yet, and it would offer an opportunity to begin to painfully reconstruct Western fabric using what we learn. Russia is not yet done. If Alexander Dugin is any indication, they have no intention of ``Americanizing'' themselves.

\hfill

\texttt{Cologero on 2011-09-17 at 19:03 said:}

Guenon always discounted the Greeks, regarding their metaphysics as less than that of India; a rather strange point of view, given that the metaphysics of Islam is based on Greek philosophy. And although in his youth he served as the secretary to the Hermetist Papus, and often referred to French Hermetists, he neglected that tradition within a Tradition. Odd again, since Hermetism played a significant role in Islam. Evola's book on Hermetism is quite insightful.

Evola was never interested in the actual content of Christianity, but rather in how the Catholic Church acted as a formative force in the Middle Ages. Since his primary goal was cementing a Nordic-Roman alliance, the church of the Eastern Empire was outside his purview.

There are two significant Orthodox writers who began as Guenonians, but later found a home in Orthodoxy. One is Seraphim Rose. The other is Phillip Sherrard who wrote a long essay on the metaphysics of Guenon in Lineaments of a Sacred Tradition. The Philokalia is churchified Platonism and probably reflects many elements and spiritual exercises from the defunct Platonic Academy. If that is where your interests lie, it is instructive to note how different it is from what is called Christianity today.

However, don't expect a mass conversion to Orthodoxy.

\hfill

\texttt{logres on 2011-09-18 at 10:29 said:}

You're probably right. Rose is very popular w/Orthodox in America, and he claimed it was Guenon who taught him how to think about religion (or to think at all metaphysically), and also rescued him from his Eastern obsession. What possibilities do you see in Western rite? If at all...

\hfill

\texttt{Charlotte Cowell on 2011-09-18 at 14:44 said:}

in ancient Greece, `Indian' referred to the Indus region, ie, the Persian empire, land of the `Indus', the river, `sindhumeans', so perhaps Guenon means this by `India'. It would make more sense in the context as Herodotus was famously rather disturbed to discover the similarities between the Greek mystery cults and the cult of Isis/Osiris found in Egypt, which in turn was informed by the Persian and Babylonian magical traditions. That said, the metaphysical system of India proper -- land of the Rishis, Maharajas, 9 wise men, Ramayana and so on -- is, I would say, far more rich and complex than what the Greeks left us (bear in mind I'm a student of the Greek school, just trying to be objective). The Greek way is much easier to grasp, especially to the western mind. Artemis and Apollo make things so simple for us, whereas the onion layers offered by Shiva, Shakti et al have relatively more facets, are more `diamond-like'.

\hfill

\texttt{logres on 2011-09-18 at 16:09 said:}

West = land of manifestations.

\hfill

\texttt{logres on 2011-09-24 at 19:33 said:}

Actually the map has a point, and is rather amusingly true. There is a lot Westerners can learn from that Church, but one has to wade through some of that to get to it. Thanks!

\hfill

\texttt{logres on 2011-09-24 at 19:35 said:}

I can't remember where I read it, but there's a legend that the East/West split is malevolent occult, whereas the North-South orientation is divinely patterned. It also happens to be oriented vertically. Is this a coincidence?


\hfill

\end{sffamily}
\end{footnotesize}

